\documentclass[11pt,a4paper]{article}
\usepackage[margin=2.5cm]{geometry}
\usepackage{amsmath, amssymb}
\usepackage{graphicx}
\usepackage{siunitx}
\usepackage{hyperref}

  \title{01\_Data\_Preparation: MF4 \texorpdfstring{$\rightarrow$}{->} Cleaned 10\,ms Dataset (Legacy Technical Note)}
\author{Validation\_Aero\_Test}
\date{\today}

\begin{document}
\maketitle

\section{Overview}
This document describes the current data-preparation pipeline implemented in
  \texttt{Code/01\_Data\_Preparation/}.

\paragraph{Note}
This file is a legacy/technical note. For the maintained documentation, see
  \texttt{Docs/01\_Data\_Preparation/01\_Data\_Preparation.tex}.

\paragraph{Goal}
Create a single analysis-ready CSV with a fixed sample raster (default: \SI{10}{ms}) and stable column names for downstream steps (run extraction and drag calculation).

\paragraph{Primary outputs}
\begin{itemize}
  \item Per-channel exports: \texttt{Outputs/01\_Data\_Preparation/Signals/*.csv}
  \item Canonical cleaned dataset: \texttt{Outputs/01\_Data\_Preparation/Aero\_Validation\_Signals\_cleaned\_10ms.csv}
\end{itemize}

\paragraph{Entry points}
\begin{itemize}
  \item \texttt{export\_channels\_10ms\_to\_signals.py}: MF4 \texorpdfstring{$\rightarrow$}{->} per-channel CSVs.
  \item \texttt{build\_cleaned\_signals\_10ms.py}: per-channel CSVs \texorpdfstring{$\rightarrow$}{->} cleaned, combined dataset with derived signals.
\end{itemize}

\section{Inputs}
\begin{itemize}
  \item Raw logging data: \texttt{Data/<session>/Stint*.mf4}
  \item Channel list: \texttt{Data/channels\_extraction.txt} (lines, \texttt{\#} for comments)
\end{itemize}

\section{Step 1: Export per-channel CSVs (fixed raster)}
Implemented in \texttt{Code/01\_Data\_Preparation/export\_channels\_10ms\_to\_signals.py}.

\subsection*{Concatenation strategy}
Each \texttt{Stint*.mf4} is treated as a separate segment:
\begin{itemize}
  \item Timestamps are shifted to start at 0 within the stint.
  \item Signals are interpolated onto a uniform grid with spacing $\Delta t$ (default: \SI{0.01}{s}).
  \item Stints are concatenated on a continuous time axis with an optional single-$\Delta t$ gap between stints (enabled by default).
\end{itemize}

\subsection*{Per-channel output schema}
For each requested channel $c$, the exporter writes:
\begin{itemize}
  \item File: \texttt{Outputs/01\_Data\_Preparation/Signals/<channel>.csv}
  \item Columns: \texttt{time} and \texttt{<channel\_name>} (numeric, NaN outside signal support)
\end{itemize}

\section{Step 2: Combine to the canonical cleaned dataset}
Implemented in \texttt{Code/01\_Data\_Preparation/build\_cleaned\_signals\_10ms.py}.

\subsection*{Join semantics}
All per-channel CSVs are read and joined on the time column using an outer join.
The resulting table is sorted by time and uses:
\begin{itemize}
  \item index name \texttt{t\_s} (seconds)
  \item output column \texttt{t\_s} after resetting the index
\end{itemize}

\subsection*{Canonical column naming}
The script renames a small set of known extracted channels to stable analysis names with units,
e.g. \texttt{SBG\_GPS1\_VELOCITY\_N} $\rightarrow$ \texttt{GPS1\_VEL\_N\_ms}.
This renaming is intentionally explicit to avoid accidental mismatches when channel naming changes.

\section{Derived signals and notation}
We use the following time-synchronized signals (after column sanitization):
\begin{itemize}
  \item $t$ [s]: timestamp (\texttt{t\_s}).
  \item $\mathbf{a}_\text{IMU} = (a_x, a_y, a_z)$ [m/s$^2$]: IMU accelerations (\texttt{IMU\_ACCEL\_X\_ms2}, \texttt{IMU\_ACCEL\_Y\_ms2}, \texttt{IMU\_ACCEL\_Z\_ms2}).
  \item $v_\text{gps}$ [m/s]: GPS speed magnitude (\texttt{GPS1\_VEL\_MAG\_ms}).
  \item $v_\text{wheel}$ [m/s]: wheelspeed magnitude, computed from \texttt{WHEEL\_SPEED\_kmh} by $v_\text{wheel} = \text{km/h} \cdot \frac{1000}{3600}$.
  \item $\theta$ [deg]: steering angle (\texttt{ICU\_SteeringAngle\_deg}).
\end{itemize}

The derived signals are produced by \texttt{aero\_validation.signals:add\_derived\_signals}.

\section{Zero-Velocity Bias Removal}
Bias and gravity components are estimated from near-stationary samples and subtracted per axis:
\begin{enumerate}
  \item Build a zero-velocity mask using available speed signals:\\
  $\mathcal{Z} = \{ i \mid v_\text{gps}[i] < 0.05 ~\text{m/s} \}$ if GPS exists; if wheelspeed exists, combine as $\mathcal{Z} = \{ i \mid v_\text{gps}[i] < 0.05 \land v_\text{wheel}[i] < 0.05 \}$.
  \item Estimate per-axis stationary bias as medians over $\mathcal{Z}$ (fallback to global median if $|\mathcal{Z}|$ is small):\\
  $b_x = \operatorname{median}\{ a_x[i] : i \in \mathcal{Z} \}$, similarly $b_y, b_z$.
  \item Form bias-corrected IMU accelerations: $a_x' = a_x - b_x$, $a_y' = a_y - b_y$, $a_z' = a_z - b_z$.
\end{enumerate}

\section{Smoothing and dv/dt}
To obtain a stable numerical derivative of speed, we apply light smoothing and then differentiate:
\begin{enumerate}
  \item Compute median sample period $\Delta t_\text{med}$ from $t$ and set a smoothing window of approximately 0.5\,s (e.g., a centered rolling mean of $\approx \frac{0.5}{\Delta t_\text{med}}$ samples).
  \item Smooth the speed trace $v(t)$ (prefer $v_\text{gps}$, fallback to $v_\text{wheel}$) to obtain $\tilde{v}(t)$.
  \item Compute the numerical derivative using a gradient operator:\\
  $$ \frac{dv}{dt}(t) \approx \operatorname{gradient}\big( \tilde{v}(t),~ t \big). $$
\end{enumerate}

\section{Straight and Moving Segment Masking}
We restrict learning and validation to segments where the vehicle is moving and approximately straight:
\begin{itemize}
  \item Moving mask: $v(t) > 0.5$\,m/s.
  \item Straight mask: $|\theta(t)| \le 5^\circ$ when steering is available; otherwise skipped.
\end{itemize}
The final validity mask is the conjunction of finiteness checks and the above conditions.

\section{Data-Driven Longitudinal Axis Estimation}
On valid samples, we estimate the longitudinal direction in the IMU frame by regressing the bias-corrected IMU signals onto $dv/dt$:
\begin{align*}
  \text{Given } \; A &= \begin{bmatrix} a_x' & a_y' & a_z' \end{bmatrix} \in \mathbb{R}^{N\times 3}, \\
  y &= \left( \frac{dv}{dt} \right) \in \mathbb{R}^{N}, \\
  \text{find } \; \mathbf{w} &\in \mathbb{R}^3 \; \text{such that } A\,\mathbf{w} \approx y \; \text{(least squares).}
\end{align*}
We solve
$$ \min_{\mathbf{w}} \; \| A\mathbf{w} - y \|_2^2, $$
and obtain $\mathbf{w}$ via a standard linear least squares solver. If insufficient valid data is available, we fall back to $\mathbf{w} = (1, 0, 0)$.

\paragraph{Longitudinal Acceleration Projection}
The longitudinal acceleration signal is then computed as
$$ a_\text{long}(t) = w_x\,a_x'(t) + w_y\,a_y'(t) + w_z\,a_z'(t). $$

\section{Notes}
\begin{itemize}
  \item \textbf{Resampling}: interpolation is performed independently per channel; NaNs indicate missing data or out-of-range timestamps.
  \item \textbf{Bias removal}: the ``stationary'' mask uses a strict low-speed threshold ($<\SI{0.05}{m/s}$) and falls back to global medians if too few stationary samples exist.
  \item \textbf{Longitudinal axis fit}: the least-squares weights are only estimated on straight + moving segments to avoid lateral acceleration contamination.
\end{itemize}

\section{Notes and Considerations}
\begin{itemize}
  \item \textbf{Smoothing}: The 0.5\,s rolling mean balances noise suppression and responsiveness before differentiation.
  \item \textbf{Bias/gravity}: Bias removal uses zero-velocity segments; grade or attitude changes can introduce offsets in $dv/dt$ vs acceleration fits.
  \item \textbf{Alignment}: If the IMU axes are slightly misaligned with the vehicle frame, the least squares method identifies an effective longitudinal combination $\mathbf{w}$.
  \item \textbf{Fallbacks}: When signals are missing or valid segments are too short, the method defaults to using the X-axis after bias removal.
  \item \textbf{Robustness}: A robust variant can remove regression outliers (e.g., via MAD-thresholding) before refitting; this is straightforward to add if needed.
\end{itemize}

\section{Implementation References}
\begin{itemize}
  \item Export: \texttt{Code/01\_Data\_Preparation/export\_channels\_10ms\_to\_signals.py}
  \item Combine + derived signals: \texttt{Code/01\_Data\_Preparation/build\_cleaned\_signals\_10ms.py}
  \item Derived signal implementation: \texttt{aero\_validation/signals.py} (function \texttt{add\_derived\_signals})
\end{itemize}

\end{document}